\usepackage[utf8]{inputenc}
\usepackage[T1]{fontenc}
% PROJECT: language-specific.  The document's main language goes last.
\usepackage[british,english,swedish]{babel}
\usepackage{booktabs}
\usepackage{array}
\usepackage{longtable}

\usepackage[
  natbib,
  style=verbose,
  citestyle=verbose,
  singletitle=false,
  maxbibnames=99,
]{biblatex}
\addbibresource{ltnotes.bib}

\usepackage[all]{foreign}
\renewcommand{\foreignfullfont}{}
\renewcommand{\foreignabbrfont}{}

\usepackage{import}
% beamer loads graphicx itself; memoir does not, and a deck with figures
% needs it in both jobs.
\usepackage{graphicx}

\usepackage[strict]{csquotes}
\usepackage[single]{acro}
% PROJECT: language-specific.  Own-language term first, the English term
% once in parentheses at first use: acro prints "long (short)" on first
% use, so the own-language word has to be the long form and the English
% expansion has to ride along inside it.
% \DeclareAcronym{IDE}{
%   short = IDE,
%   long  = {utvecklingsmiljö
%            (\foreignlanguage{english}{integrated development
%             environment})},
% }

\usepackage{subcaption}

\usepackage{xparse}

\let\email\texttt

% minted v3+ (TeX Live 2024+) auto-detects the output directory; the old
% [outputdir=...] package option is an error there.
\usepackage{minted}
\setminted{autogobble,fontsize=\footnotesize}

% PythonTeX runs the example programs and embeds their real terminal
% sessions (didactic's \runpython[transcript]), so the deck never
% hand-copies its own output: pythontex runs from ltxobj (the outputdir),
% and workingdir=.. moves it up to the deck directory so examples/...
% resolves.
\usepackage[makestderr]{pythontex}
\setpythontexoutputdir{.}
\setpythontexworkingdir{..}

\usepackage{amsmath}
\usepackage{amssymb}
\usepackage{mathtools}
\usepackage{amsthm}
\usepackage{thmtools,thm-restate}
\usepackage[unq]{unique}

\usepackage[binary-units,parse-numbers=false]{siunitx}

\makeatletter
\@ifclassloaded{beamer}
  {\relax}
  {\usepackage[colorlinks]{hyperref}}
\makeatother
% PROJECT: language-specific.  cleveref reads its language from its own
% package options / global class options, not from babel's package
% options, so name the language here too.  No [capitalize]: Swedish
% typography wants lowercase mid-sentence references (use \Cref for a
% sentence-initial one).
\usepackage[swedish]{cleveref}
% cleveref's Swedish names call the appendix type "Appendix", but memoir
% and babel call it "Bilaga".  Align cleveref with them, in an
% \AtBeginDocument hook registered after cleveref's own (from the language
% option) so ours wins.
\AtBeginDocument{%
  \crefname{appendix}{bilaga}{bilagor}%
  \Crefname{appendix}{Bilaga}{Bilagor}%
  % The prose refers back to numbered lines of a program listing as
  % "rader", not cleveref's Swedish default "punkt".
  \crefname{enumi}{rad}{raderna}%
  \Crefname{enumi}{Rad}{Raderna}%
  \crefname{footnote}{fotnot}{fotnoterna}%
  \Crefname{footnote}{Fotnot}{Fotnoterna}%
}

% footnotereuse=compact: the appendices cite the same sources many times,
% often several times on one page.  didactic lets a repeat citation of the
% same source (same key, prenote and postnote) on the same page reuse the
% earlier margin note instead of adding one; the default, fast, gives up
% reuse and reserves a blank as soon as a citation's page changes between
% passes, compact only when it changes back again.  That leaves fewer
% notes and blanks in the margin, at the cost of more passes.
\usepackage[marginparmargin=outer,footnotereuse=compact]{didactic}
% didactic sets footnotes --- and, under the verbose citation style, every
% \autocite --- as margin notes, and LaTeX allocates \marginpar boxes from
% the same pool of eighteen boxes as floats.  A chapter with this many
% margin references empties the pool and dies with "Too many unprocessed
% floats" pages later, in a place that has nothing to do with the cause.
\extrafloats{200}
% memoir places a side caption with \rlap, so it takes part in no
% collision avoidance at all and silently overprints an \ltnote that
% happens to sit at the same height.  Aligning the caption with the BOTTOM
% of the figure moves it below the notes, which anchor where they stand in
% the body text.  (\setsidecappos is memoir's; beamer does not define it.)
\ifdefined\setsidecappos\setsidecappos{b}\fi
% The form of a construct --- its lines, indentation, colons and
% parentheses, when it is used and what usually comes with it --- in one
% environment of its own, so that every such box in the deck looks the
% same: "Syntax: if" on the slides, "Syntax 3 (if)" in the notes.  It is
% didactic's own mechanism for a semantic environment, like lo below.
\ProvideSemanticEnv{syntax}[block]{Syntax}
  [style=definition,numbered=yes]
  {syntax}{syntaxes}
  {Syntax}{Syntaxes}
\ProvideTranslation{swedish}{Syntax}{Syntax}
\ProvideTranslation{swedish}{syntax}{syntax}
\ProvideTranslation{swedish}{syntaxes}{syntaxer}
\ProvideTranslation{swedish}{Syntaxes}{Syntaxer}
\ProvideTranslation{british}{Syntax}{Syntax}
\ProvideTranslation{british}{syntax}{syntax}
\ProvideTranslation{british}{syntaxes}{syntaxes}
\ProvideTranslation{british}{Syntaxes}{Syntaxes}

% Pictures of Python's memory model and control flow, shared by every deck
% (pythonfigs.sty is a link to modules/pythonfigs.sty).
\usepackage{pythonfigs}

\ProvideSemanticEnv{lo}{Learning Objective}
  [style=definition,numbered=yes]
  {LO}{LO}
  {Learning objective}{Learning objectives}
\ProvideTranslation{british}{Learning Objective}{Learning Objective}
\ProvideTranslation{british}{LO}{LO}
\ProvideTranslation{british}{Learning objective}{Learning objective}
\ProvideTranslation{british}{Learning objectives}{Learning objectives}
% PROJECT: language-specific.
\ProvideTranslation{swedish}{Learning Objective}{Lärandemål}
\ProvideTranslation{swedish}{LO}{lm}
\ProvideTranslation{swedish}{Learning objective}{Lärandemål}
\ProvideTranslation{swedish}{Learning objectives}{Lärandemål}

% Paper titles in the generated hit-list tables can carry Greek letters;
% scholar's table export sets them in math mode itself, and this mapping
% stays as a guard for the same characters typed anywhere else.
\DeclareUnicodeCharacter{03C0}{\ensuremath{\pi}}
\DeclareUnicodeCharacter{03BB}{\ensuremath{\lambda}}
\DeclareUnicodeCharacter{03B1}{\ensuremath{\alpha}}
\DeclareUnicodeCharacter{03B2}{\ensuremath{\beta}}
\DeclareUnicodeCharacter{03BC}{\ensuremath{\mu}}
